% Generated from THESIS_DRAFT.md.
\chapter*{Tóm tắt}
\thispagestyle{empty}

Trong thao tác robot chuỗi dài, các quan sát hiện tại có thể không cho biết một chu kỳ đã hoàn tất hay chưa, dù quyết định tiếp theo phụ thuộc vào lịch sử tương tác. Khóa luận xây dựng bộ nhớ sự kiện có bằng chứng cho mô hình thị giác - ngôn ngữ - hành động (VLA), nhằm lưu giữ các mốc tiến độ đã được xác nhận và cung cấp chúng qua một bản tóm tắt ngôn ngữ gọn nhẹ phù hợp với giới hạn ngữ cảnh.

Mỗi sự kiện được gắn với bằng chứng quan sát và thời điểm khả dụng; khi lịch sử thiếu hoặc bị gián đoạn, hệ thống giữ trạng thái chưa xác định thay vì suy diễn tiến độ. Khóa luận chứng minh hai tính chất hình thức trong phạm vi các giả định đã nêu: bộ nhớ sự kiện chỉ thay đổi khi có chuyển tiếp hợp lệ, và nội dung cung cấp cho chính sách tại thời điểm hiện tại chỉ phụ thuộc dữ liệu đã có đến thời điểm đó. Bản ghi được chiếu thành văn bản và dùng để thích ứng các chính sách VLA.

Đánh giá thực nghiệm của nghiên cứu bao gồm phân tích tiền huấn luyện, can thiệp vào prompt chỉ dẫn hành động, đo đạc điều khiển vòng kín, kiểm chứng cơ chế dừng ngoại vi dựa trên bộ nhớ và đo mức duy trì hiệu năng trên tác vụ gốc sau tinh chỉnh. Kết quả xác nhận cơ chế truy vấn theo sự kiện vận hành được và cho thấy mô hình phản ứng với nội dung bộ nhớ trong các trường hợp khảo sát. Dù lợi ích trực tiếp qua lời nhắc còn hạn chế đối với tỷ lệ thành công nội tại, việc khai thác bằng chứng bộ nhớ cho cơ chế dừng ngoại vi bước đầu giúp chuyển hóa các chu kỳ thao tác hoàn thành thành thành công toàn phần. Mức duy trì hiệu năng sau tinh chỉnh cũng khác biệt rõ rệt giữa hai mô hình nền. Các kết quả này làm rõ khả năng vận hành và giới hạn thực nghiệm của phương pháp, đồng thời đặt ra nhu cầu xác nhận trên tập trạng thái rộng hơn với nhiều ngẫu nhiên hơn.

\textbf{Từ khóa:} Mô hình thị giác - ngôn ngữ - hành động, Tác vụ long-horizon, Bộ nhớ sự kiện, Quan sát một phần, Thao tác robot.

